\documentclass[12pt, a4paper]{article}

% Pakiety językowe i kodowanie
\usepackage[utf8]{inputenc}
\usepackage[polish]{babel}
\usepackage[T1]{fontenc}

% Pakiety do formatowania strony i tabel
\usepackage[margin=2.5cm]{geometry}
\usepackage{array}
\usepackage{multicol} % Wymagany do scalania kolumn

% Pakiet do customizacji list i numeracji
\usepackage{enumitem}

\title{Sprawozdanie z Laboratorium}
\author{Filip Będkowski, Danylo Drygo}
\date{\today}

\begin{document}

\noindent
\begin{tabular}{|p{0.48\textwidth}|p{0.48\textwidth}|}
    \hline
    \multicolumn{2}{|l|}{\textit{Lab4 Omówienie projektów}} \\ \hline
    \textit{Filip Będkowski, Daniel Cebula} & {'Okazjomat'} \\ \hline
\end{tabular}

\vspace{0.8cm}

\section*{Co zostało wprowadzone}

% Ustawienie formatowania dla wielopoziomowej listy
\begin{enumerate}[label=\arabic*.]
    \item \textbf{Określona struktura programu}
    \begin{enumerate}
        \item Podział programu na Frontend oraz Backend
        \item \textbf{Frontend} realizowany jest przez framework React+Vite+JS 
        \item \textbf{Backend} realizowany jest w Python
        \item w pliku .gitignore zaznaczone są pliki pomocnicze i konfiguracyjne edytora vsc, logów, plików tymczasowych
    \end{enumerate}

    \item \textbf{Użyte biblioteki}
    \begin{enumerate}[label=\arabic{enumi}.\arabic*.]
        \item \textbf{Frontend}
        \begin{itemize}[label=--]
            \item React - biblioteka ui
            \item Inline / Embedded SVG - system ikon
            \item CSS / CSS Modules - stylowanie
        \end{itemize}
        \item \textbf{Backend}
        \begin{itemize}[label=--]
            \item Pandas - obsługa baz danych
            \item Crawlee[playwright] - pobieranie danych HTML z obsługą JS
            \item \textsl{biblioteki w pliku requirements.txt}
            \item FastAPI do integracji Frontendu z Backendem
        \end{itemize}
    \end{enumerate}

    \item \textbf{Standardy kodowania} 
            \begin{center}
                Do kodowania używamy programu \textsl{Visual Studio Code} z rozszerzeniami \textsl{prettier} do HTML, CSS, JS, a do Pythona używamy \textsl{Pylance}  
            \end{center}
    
    \item \textbf{Co dalej}
            \begin{itemize}
                \item Wprowadzenie pełnej integracji między Frontend i Backend
                \item Zmiana z Pythona na Rusta ze względu na prędkość pobierania danych
                \item stworzenie prostego algorytmu analizy cen
                \item dodanie możliwości tworzenia własnego algorytmu dotyczącego jednego obszaru rynku 
            \end{itemize}
\end{enumerate}

\end{document}